% from ega1-6-fr.tex
\oldpage[I]{148}
un idéal maximal de $A$, autrement dit tels que $A/\mathfrak{j}_x$ soit un
corps (nécessairement égal à $k(x)$). Comme $A/\mathfrak{j}_x$ est une
$K$-algèbre de type fini, on voit que si $x$ est fermé, $k(x)$ est un corps
qui est une algèbre de type fini sur $K$, donc nécessairement une $K$-algèbre
de \emph{rang fini} \cite{I-21}. Inversement, si $k(x)$ est de rang fini sur
$K$, il en est de même de $A/\mathfrak{j}_x\subset k(x)$ et comme tout anneau
intègre qui est une $K$-algèbre de rang fini est un corps, on a
$A/\mathfrak{j}_x=k(x)$, donc $x$ est fermé.

\begin{corollary}[6.4.3]
\label{I.6.4.3-fr}
Soient $K$ un corps algébriquement clos, $X$ un $K$-préschéma algébrique~; les
points fermés de $X$ sont alors les points rationnels sur $K$
(\hyperref[I.3.4.4-fr]{3.4.4}) et s'identifient canoniquement aux points de
$X$ à valeurs dans $K$.
\end{corollary}

\begin{proposition}[6.4.4]
\label{I.6.4.4-fr}
Soit $X$ un préschéma algébrique sur un corps $K$. Les propriétés suivantes
sont équivalentes~:
\begin{enumerate}
  \item[\rm(a)] $X$ est artinien.
  \item[\rm(b)] L'espace sous-jacent à $X$ est discret.
  \item[\rm(c)] L'espace sous-jacent à $X$ n'a qu'un nombre fini de points
    fermés.
  \item[\rm(c')] L'espace sous-jacent à $X$ est fini.
  \item[\rm(d)] Les points de $X$ sont fermés.
  \item[\rm(e)] $X$ est isomorphe à $\operatorname{Spec}(A)$, où $A$ est une
    $K$-algèbre de rang fini.
\end{enumerate}
\end{proposition}

Comme $X$ est noethérien, il résulte de (\hyperref[I.6.2.2-fr]{6.2.2}) que
les conditions a), b), d) sont équivalentes et entraînent c) et c')~; par
ailleurs, il est clair que e) entraîne a). Reste à voir que c) entraîne d) et
e)~; on peut se borner au cas où $X$ est affine. Alors $A(X)$ est une
$K$-algèbre de type fini (\hyperref[I.6.3.3-fr]{6.3.3}), donc un anneau de
Jacobson (\cite[p.~3-11 et 3-12]{I-1}), dans lequel il n'y a par hypothèse
qu'un nombre fini d'idéaux maximaux. Comme une intersection finie d'idéaux
premiers ne peut être un idéal premier que si elle est égale à l'un d'eux,
tout idéal premier de $A(X)$ est donc maximal, d'où d). En outre, on sait
alors (\hyperref[I.6.2.2-fr]{6.2.2}) que $A(X)$ est une $K$-algèbre
artinienne de type fini, donc nécessairement de \emph{rang fini} \cite{I-21}.

\begin{env}[6.4.5]
\label{I.6.4.5-fr}
Lorsque les conditions de (\hyperref[I.6.4.4-fr]{6.4.4}) sont satisfaites,
on dit que $X$ est un schéma \emph{fini sur $K$}
(\emph{cf.} (\hyperref[II.6.1.1-fr]{II, 6.1.1})), ou un
\emph{$K$-schéma fini}, de \emph{rang} $[A:K]$ que l'on note aussi
$\operatorname{rg}_K(X)$~; si $X$, $Y$ sont deux schémas finis sur $K$, on a
\[
  \operatorname{rg}_K(X\amalg Y)
  =\operatorname{rg}_K(X)+\operatorname{rg}_K(Y)
  \tag{6.4.5.1}\label{I.6.4.5.1-fr}
\]
\[
  \operatorname{rg}_K(X\times_K Y)
  =\operatorname{rg}_K(X)\operatorname{rg}_K(Y)
  \tag{6.4.5.2}\label{I.6.4.5.2-fr}
\]
comme il résulte de (\hyperref[I.3.2.2-fr]{3.2.2}).
\end{env}

\begin{corollary}[6.4.6]
\label{I.6.4.6-fr}
Soit $X$ un schéma fini sur un corps $K$. Pour toute extension $K'$ de $K$,
$X\otimes_K K'$ est un schéma fini sur $K'$, et son rang sur $K'$ est égal au
rang de $X$ sur $K$.
\end{corollary}

En effet, si $A=A(X)$, on a $[A\otimes_K K':K']=[A:K]$.

\begin{corollary}[6.4.7]
\label{I.6.4.7-fr}
Soit $X$ un schéma fini sur un corps $K$~; on pose
$n=\sum_{x\in X}[k(x):K]_s$
\emph{(on rappelle que si $K'$ est une extension de $K$, $[K':K]_s$ est le
\emph{rang séparable} de $K'$ sur $K$, rang de la plus grande extension
algébrique séparable de $K$ contenue dans $K'$)}~;
